\section{Bases de Diseño}

\subsection{Condiciones Ambientales del Sitio}

Las condiciones ambientales del sitio se presentan en la Tabla \ref{tab:environmental_conditions}. Estos valores corresponden a las condiciones del sitio de instalación en Cajicá, Cundinamarca (planta BRINSA) \cite{dml_informe_investigacion_2026} y se utilizan para la estimación de las propiedades del aire de impulsión.

\begin{table}[htbp]
	\centering
	\caption{Condiciones ambientales del sitio de instalación}
	\label{tab:environmental_conditions}
	\setcounter{itemcount}{0}
	\begin{tabularx}{\textwidth}{@{}NXclc@{}}
		\toprule
		\multicolumn{1}{c}{\textbf{Item}} & \textbf{Parámetro Ambiental} & \textbf{Valor} & \textbf{Unidad} & \textbf{Referencia} \\
		\midrule
		& Temperatura ambiente máxima & 21.0 & $^{\circ}$C & IDEAM 2023 \\
		& Temperatura ambiente mínima & 3.0 & $^{\circ}$C & IDEAM 2023 \\
		& Temperatura ambiente promedio & 14.0 & $^{\circ}$C & IDEAM 2023 \\
		& Humedad relativa promedio & 84.0 & \% & IDEAM 2023 \\
		& Presión atmosférica & 74.1 & kPa & NOAA 2022 \\
		& Altitud sobre el nivel del mar & 2 558 & msnm & IGAC 2021 \\
		\bottomrule
	\end{tabularx}
\end{table}

\subsection{Propiedades del Aire de Impulsión}

El aire de impulsión se considera a 21 $^{\circ}$C (condición de diseño de verano 0.4 \%) y 74.1 kPa (condiciones del sitio), con propiedades aproximadas presentadas en la Tabla \ref{tab:air_properties}.

\begin{table}[htbp]
	\centering
	\caption{Propiedades del aire de impulsión}
	\label{tab:air_properties}
	\setcounter{itemcount}{0}
	\begin{tabularx}{\textwidth}{@{}NXclc@{}}
		\toprule
		\multicolumn{1}{c}{\textbf{Item}} & \textbf{Propiedad} & \textbf{Valor} & \textbf{Unidad} & \textbf{Referencia} \\
		\midrule
		& Densidad & 0.88 & kg/m$^3$ & Gas ideal a condiciones del sitio \\
		& Viscosidad dinámica & $1.81 \times 10^{-5}$ & Pa$\cdot$s & Cengel \& Ghajar (2020) \\
		& Temperatura de impulsión & 14.0 & $^{\circ}$C & Criterio de diseño \\
		& Calidad de filtración mínima & MERV 13-14 & -- & ASHRAE 52.2 \\
		\bottomrule
	\end{tabularx}
\end{table}

\subsection{Datos del Recinto y Criterios de Ventilación}

La Tabla \ref{tab:room_data} resume las características geométricas y los criterios de ventilación del laboratorio. El laboratorio se ubica dentro de una planta de hipoclorito de calcio, cuyo ambiente exterior es altamente corrosivo; el ventilador, las rejillas y los soportes expuestos deben ser en PRFV, acero inoxidable 316 o acero con recubrimiento epóxico, y el motor debe ser TEFC anticorrosivo.

\begin{table}[htbp]
	\centering
	\caption{Datos del recinto y criterios de ventilación}
	\label{tab:room_data}
	\setcounter{itemcount}{0}
	\begin{tabularx}{\textwidth}{@{}N>{\hsize=0.9\hsize}X>{\hsize=1.2\hsize}Xc>{\hsize=0.9\hsize}X@{}}
		\toprule
		\multicolumn{1}{c}{\textbf{Item}} & \textbf{Parámetro} & \textbf{Valor} & \textbf{Unidad} & \textbf{Referencia} \\
		\midrule
		& Volumen efectivo del laboratorio & 320 & m$^3$ & Medición / modelo 3D \\
		& Renovaciones de aire requeridas & 12 & ren/h & ASHRAE 170 / WHO / NIH \\
		& Caudal de diseño & 3 840 (64.0 m$^3$/min; 2 260 CFM) & m$^3$/h & Memoria de cálculo \\
		& Estrategia de ventilación & Ventilador directo + rejillas de descarga libre & -- & Criterio de diseño \\
		& Descarga de rejillas & Libre a la atmósfera (sin contrapresión) & -- & Criterio de diseño \\
		\bottomrule
	\end{tabularx}
\end{table}

\subsection{Criterios de Diseño del Ventilador}

La Tabla \ref{tab:fan_criteria} presenta los criterios de diseño del ventilador de impulsión, de tipo axial mural (placa mural) Ø560 mm con transmisión directa, configuración uniforme con el montaje típico instalado en la planta (cubierta intemperie con banco de filtración, estructura de unión pernada al muro y malla de protección interior); el motor va en el cubo del impulsor dentro de la corriente de aire, en ejecución encapsulada severe duty anticorrosiva. El punto de trabajo en el sitio corresponde al escenario de filtración cargada ($\Delta P$ filtro 154 Pa + 11 Pa de pérdida en las rejillas de descarga libre = 165 Pa total); dado que las curvas de catálogo se publican a densidad estándar ($\rho$ = 1.2 kg/m$^3$), la selección comercial se realiza en el punto equivalente de catálogo 3 840 m$^3$/h a 225 Pa \cite{dml_informe_investigacion_2026,dml_hd_vent_001}.

\begin{table}[htbp]
	\centering
	\caption{Criterios de diseño del ventilador de impulsión}
	\label{tab:fan_criteria}
	\setcounter{itemcount}{0}
	\begin{tabularx}{\textwidth}{@{}N>{\hsize=0.9\hsize}X>{\hsize=1.2\hsize}Xc>{\hsize=0.9\hsize}X@{}}
		\toprule
		\multicolumn{1}{c}{\textbf{Item}} & \textbf{Parámetro} & \textbf{Valor} & \textbf{Unidad} & \textbf{Referencia} \\
		\midrule
		& Tipo de ventilador & Axial mural (placa mural) Ø560 mm PRFV/epóxico, transmisión directa (uniformidad con montaje típico de planta) & -- & Criterio de diseño \\
		& Caudal de diseño & 3 840 (1.0667 m$^3$/s) & m$^3$/h & Memoria de cálculo \\
		& Presión total en el sitio (filtro cargado) & 165 & Pa & Memoria de cálculo \\
		& Punto equivalente de catálogo ($\rho$ = 1.2 kg/m$^3$) & 225 & Pa & Corrección de densidad, $k$ = 0.733 \\
		& Presión total en el sitio (filtro limpio) & 70 & Pa & Memoria de cálculo \\
		& Punto equivalente de catálogo (limpio) & 95 & Pa & Corrección de densidad, $k$ = 0.733 \\
		& Eficiencia total supuesta & 0.55 & -- & Criterio de diseño (provisional) \\
		& Potencia teórica de aire & 0.320 (0.43 HP) & kW & Ec. (\ref{eq:potencia}) \\
		& Motor provisional & 0.75 HP TEFC encapsulado anticorrosivo (en la corriente de aire), 440 V, 3$\phi$, 60 Hz (margen de servicio 1.5, a confirmar con catálogo) & -- & Criterio de diseño / cliente \\
		& Diámetro del impulsor (seleccionado) & 560 & mm & Uniformidad montaje típico de planta (REV2) \\
		& Velocidad real en boca (Ø560 mm) & 4.33 & m/s & $Q$ / $A_\text{boca}$ \\
		& Velocidad de impulsión (BC del CFD, histórica) & 8.0 & m/s & Criterio de diseño CFD \\
		& Área de boca de impulsión (BC del CFD, histórica) & 0.1333 & m$^2$ & Ec. (\ref{eq:area_vent}) \\
		& Diámetro equivalente (radio) (BC del CFD, histórica) & 412 (206.0) & mm & Ec. (\ref{eq:diametro}) \\
		\bottomrule
	\end{tabularx}
\end{table}

\subsection{Escenarios de Filtración}

La calidad de filtración definitiva es MERV 13-14: el laboratorio es de análisis industrial y no de biocontención, por lo que la filtración HEPA quedó evaluada y descartada. La Tabla \ref{tab:filtration_scenarios} presenta los escenarios de filtración considerados en el diseño, con las caídas de presión corregidas a la densidad del sitio.

\begin{table}[htbp]
	\centering
	\caption{Escenarios de filtración del sistema de impulsión}
	\label{tab:filtration_scenarios}
	\setcounter{itemcount}{0}
	\begin{tabularx}{\textwidth}{@{}NXcccc@{}}
		\toprule
		\multicolumn{1}{c}{\textbf{Item}} & \textbf{Escenario} & \textbf{$\Delta P$ Filtro} & \textbf{$\Delta P$ Total} & \textbf{Potencia} & \textbf{Estado} \\
		& & \textbf{(Pa)} & \textbf{(Pa)} & \textbf{(kW)} & \\
		\midrule
		& MERV 13-14 limpio & 59 & 70 & 0.136 & Operación inicial \\
		& MERV 13-14 cargado (diseño) & 154 & 165 & 0.320 & Punto de selección \\
		& HEPA H13 limpio & 250 & 261 & 0.506 & Referencia histórica (no aplica) \\
		& HEPA H13 cargado & 600 & 611 & 1.185 & Referencia histórica (no aplica) \\
		\bottomrule
	\end{tabularx}
\end{table}

\subsection{Rejillas de Exfiltración}

La Tabla \ref{tab:grille_criteria} presenta los criterios de diseño de las rejillas de descarga. Las rejillas descargan libremente a la atmósfera: a la densidad del sitio, el cierre de orificio ($C_d$ = 0.60) de las tres rejillas a 3.0 m/s representa una pérdida de descarga de 11 Pa, sin set-point de presión en el recinto \cite{dml_hd_rej_001}.

\begin{table}[htbp]
	\centering
	\caption{Criterios de diseño de las rejillas de exfiltración}
	\label{tab:grille_criteria}
	\setcounter{itemcount}{0}
	\begin{tabularx}{\textwidth}{@{}N>{\hsize=0.9\hsize}X>{\hsize=1.2\hsize}Xc>{\hsize=0.9\hsize}X@{}}
		\toprule
		\multicolumn{1}{c}{\textbf{Item}} & \textbf{Parámetro} & \textbf{Valor} & \textbf{Unidad} & \textbf{Referencia} \\
		\midrule
		& Número de rejillas & 3 & unidades & Criterio de diseño \\
		& Dimensiones por rejilla & 353 $\times$ 336 & mm & Memoria de cálculo \\
		& Velocidad facial & 3.0 & m/s & Criterio de diseño \\
		& Área neta total (por rejilla) & 0.3556 (0.1185) & m$^2$ & Ec. (\ref{eq:area_exfil}) \\
		& Presión de descarga libre ($C_d$ = 0.60) & 11 & Pa & Cierre de orificio, $\rho$ = 0.88 kg/m$^3$ \\
		& Malla anti-insectos & Inox 316, 18$\times$18 (abertura $\approx$1 mm) & -- & Criterio de diseño \\
		\bottomrule
	\end{tabularx}
\end{table}

\subsection{Corrosividad del Ambiente y Criterio de Materiales}

El laboratorio se ubica dentro de una planta de hipoclorito de calcio, cuyo ambiente exterior es altamente corrosivo (cloro, hipoclorito y humedad alta). La Tabla \ref{tab:material_criteria} resume el criterio de materiales adoptado para los componentes expuestos \cite{dml_informe_investigacion_2026}.

\begin{table}[htbp]
	\centering
	\caption{Criterio de materiales para el ambiente corrosivo del sitio}
	\label{tab:material_criteria}
	\setcounter{itemcount}{0}
	\begin{tabularx}{\textwidth}{@{}NlX@{}}
		\toprule
		\multicolumn{1}{c}{\textbf{Item}} & \textbf{Elemento} & \textbf{Criterio de Material} \\
		\midrule
		& Ventilador (carcasa y rueda) & PRFV con resina viniléster (primera opción) o polipropileno \\
		& Rejillas, malla anti-insectos y ferretería & Acero inoxidable 316/316L (clase A4) \\
		& Soportes expuestos & PRFV, inox 316 o acero con recubrimiento epóxico (segunda línea, con inspección programada) \\
		& Motor & TEFC anticorrosivo, ejecución severe duty \\
		& Acero galvanizado & Descartado en la corriente y a la intemperie \\
		\bottomrule
	\end{tabularx}
\end{table}

\subsection{Normativas Aplicables}

El número de 12 renovaciones de aire por hora se sustenta en los rangos recomendados por las normas y guías internacionales presentadas en la Tabla \ref{tab:applicable_norms}.

\begin{table}[htbp]
	\centering
	\caption{Normativas y guías aplicables al diseño de ventilación}
	\label{tab:applicable_norms}
	\setcounter{itemcount}{0}
	\begin{tabularx}{\textwidth}{@{}N>{\hsize=1.2\hsize}X>{\hsize=1.0\hsize}X>{\hsize=0.8\hsize}X@{}}
		\toprule
		\multicolumn{1}{c}{\textbf{Item}} & \textbf{Norma / Guía} & \textbf{Recomendación ACH} & \textbf{Referencia} \\
		\midrule
		& ASHRAE 170 — Ventilation of Health Care Facilities & 6–15 ACH según uso & ASHRAE (2021) \\
		& ASHRAE 62.1 — Ventilation for Acceptable IAQ & Caudal por ocupación y área & ASHRAE (2022) \\
		& NFPA 99 — Health Care Facilities Code & Criterios de ventilación de espacios de salud & NFPA (2021) \cite{nfpa99_2021} \\
		& WHO Laboratory Biosafety Manual (BSL-2) & 6–12 ACH & WHO (2004) \\
		& WHO Laboratory Biosafety Manual (BSL-3) & 12–15 ACH & WHO (2004) \\
		& NIH Design Requirements Manual & 10–12 ACH & NIH (2023) \\
		& OSHA 29 CFR 1910.1450 & Control de exposición por ventilación & OSHA (2012) \\
		& ASHRAE Handbook Fundamentals & Coeficiente de descarga de orificios ($C_d$ = 0.60) & ASHRAE (2021) \\
		& RETIE & Suministro y protección del motor & MinCIT (vigente) \\
		& NTC 2050 — Código Eléctrico Colombiano & Circuitos y protecciones & ICONTEC (vigente) \\
		& Resolución 0312 de 2019 & Habilitación de servicios de salud (si aplica) & MinSalud (2019) \\
		\bottomrule
	\end{tabularx}
\end{table}

La selección de 12 ACH cubre con holgura el rango superior establecido para laboratorios BSL-2 y el umbral inferior para laboratorios BSL-3, proporcionando un margen de seguridad respecto al mínimo de 6 ACH para espacios generales.
